\clearpage\thispagestyle{empty}
\pdfbookmark[0]{Quellen zur Anzeige der Theoria motus von 1809}{anzeige1809-sources}
\begin{center}\textbf{Quellen zur Anzeige der \textit{Theoria motus} von 1809}\end{center}
\noindent Diese KI-Übertragung des deutschen Originals in heutiges Deutsch folgt dem Text in Gauß’ \textit{Werken}, Band~VI (Göttingen, 1874), S.~51 und 53--60. Sie umfasst das Abteilungstitelblatt und die Anzeige aus den \textit{Göttingischen gelehrten Anzeigen}, 96.~Stück, 17.~Juni 1809, S.~945--955. \href{https://www.digitale-sammlungen.de/de/view/bsb10538675?page=947}{Erstdruck in der Bayerischen Staatsbibliothek}. \href{https://gdz.sub.uni-goettingen.de/id/PPN236007467}{Historischer deutscher Text in den Werken}.
\par\smallskip\noindent Die Vorrede zum Buch ist ein eigener Text. Sie liegt in Charles Henry Davis’ englischer Übersetzung der \textit{Theoria motus} vor (Boston, 1857). Zur heutigen Terminologie der Bahnbestimmung siehe Ernst Friedrich Maria Jochim, \textit{Satellitenbewegung}, Band~III, Fassung vom 9.~Februar 2022, S.~342 und 363--364.
\par\smallskip\noindent Übertragung und Satz wurden mit KI erstellt.
\clearpage
